\begin{trapbox}
	\begin{description}[style=nextline, leftmargin=2.8em, labelsep=0.5em, itemsep=0.35em]
		\item[\textbf{\textcolor{CTrap}{易错点一（量词未换）}}] 将 $\neg(\forall x\in M,\ p(x))$ 错误地写成 $\forall x\in M,\ \neg p(x)$。
		\item[\textbf{\textcolor{CTrap}{正确理解（标准形式）：}}] 必须同时改变量词，即 $\exists x\in M,\ \neg p(x)$。
		\item[\textbf{\textcolor{CTrap}{错因分析（记忆混淆）：}}] 记错了规则，误以为“所有不满足”就是原命题的否定。
	\end{description}

	\medskip
	\begin{description}[style=nextline, leftmargin=2.8em, labelsep=0.5em, itemsep=0.35em]
		\item[\textbf{\textcolor{CTrap}{易错点二（谓词未否）}}] 将 $\neg(\forall x\in M,\ p(x))$ 错误地写成 $\exists x\in M,\ p(x)$（只改量词，不否谓词）。
		\item[\textbf{\textcolor{CTrap}{正确理解（标准规则）：}}] 必须同时否定谓词，即 $\exists x\in M,\ \neg p(x)$。
		\item[\textbf{\textcolor{CTrap}{错因分析（粗心省略）：}}] 写否定时只记得换量词，忘记把条件取反。
	\end{description}

	\medskip
	\begin{description}[style=nextline, leftmargin=2.8em, labelsep=0.5em, itemsep=0.35em]
		\item[\textbf{\textcolor{CTrap}{易错点三（论域消失）}}] 将否定写成 $\exists x,\ \neg p(x)$，省略了“$\in M$”。
		\item[\textbf{\textcolor{CTrap}{正确理解（完整表达）：}}] 量词必须带上论域，即 $\exists x\in M,\ \neg p(x)$。
		\item[\textbf{\textcolor{CTrap}{错因分析（忽视结构）：}}] 认为量词已经隐含范围，但严格逻辑要求必须明确论域。
	\end{description}
\end{trapbox}
